\documentclass[11pt,a4paper]{article}

\usepackage[margin=2.2cm]{geometry}
\usepackage[T1]{fontenc}
\usepackage{lmodern}
\usepackage{microtype}
\usepackage{booktabs}
\usepackage{tabularx}
\usepackage{array}
\usepackage{hyperref}
\usepackage{xcolor}
\usepackage{listings}
\usepackage{enumitem}
\usepackage{parskip}

\hypersetup{
  colorlinks=true,
  linkcolor=black,
  urlcolor=blue,
  pdftitle={Job fetching: from trigger to Postgres},
}

\lstset{
  basicstyle=\ttfamily\footnotesize,
  breaklines=true,
  columns=fullflexible,
  frame=single,
  rulecolor=\color{black!20},
  backgroundcolor=\color{black!3},
  xleftmargin=0.5em,
  xrightmargin=0.5em,
}
\setlist[itemize]{leftmargin=*,itemsep=0.2em}
\setlist[enumerate]{leftmargin=*,itemsep=0.2em}

\title{Job fetching in Kuchup:\\from trigger to Postgres}
\author{Verified against repository code}
\date{14 September 2026}

\begin{document}
\maketitle

\begin{abstract}
This document describes how Kuchup \textbf{fetches jobs}: how a run starts,
how each company's ATS or career page is scraped, how listings are filtered
and merged, and how rows land in Postgres. Source of truth is the code on
\texttt{main}. Fetch writes the shared catalog
(\texttt{companies}, \texttt{matching\_jobs}, \texttt{country\_meta}) plus
operational tables (\texttt{fetch\_runs}, \texttt{company\_fetch\_attempts}).
It does \textbf{not} write per-user tracking.
There is no separate ATS-cache table: cached ATS identity lives on
\texttt{companies.ats\_type} / \texttt{companies.ats\_url}.
\end{abstract}

\tableofcontents
\newpage

\section{Purpose: fetching jobs vs company discovery}

\textbf{Fetching jobs} means: take companies that already exist in the catalog,
hit their careers URL or ATS API, keep software-relevant listings, merge them
into \texttt{matching\_jobs}, and record a fetch run.

\textbf{Company discovery} (\texttt{build\_companies}) is a different pipeline.
It does not scrape openings.

\begin{center}
\begin{tabularx}{\textwidth}{l X X}
  \toprule
  \textbf{Concern} & \textbf{Discovery} & \textbf{Fetch} \\
  \midrule
  Entry &
    \texttt{scripts/build\_companies.py} &
    \texttt{apps/fetch-worker/run.py} or panel Fetch \\
  Input &
    Country key (optional company name) &
    Existing \texttt{companies} rows \\
  Work &
    Find \texttt{careers\_url} from relocate.me &
    Detect ATS, list jobs, filter, merge, enrich \\
  Writes &
    Company metadata &
    Jobs, fetch flags, run/attempt rows \\
  \bottomrule
\end{tabularx}
\end{center}

README architecture:

\begin{lstlisting}
relocate.me -> build_companies -> Postgres catalog
            -> scrape / v2 fetch -> Postgres catalog
\end{lstlisting}

After a successful country or company \textbf{run},
\texttt{fetch/runner.py} enqueues
\texttt{enqueue\_country\_opportunity\_refresh} so the Go role-propagator can
rebuild per-user opportunity slots. That is a side effect after catalog write,
not part of scrape itself.

\section{Triggers}

A fetch can start in four ways. Only one run may be
\texttt{fetch\_runs.status = running} at a time (global lock across panel and
worker).

\subsection{Scheduled worker (production path)}

\begin{center}
\begin{tabularx}{\textwidth}{l X}
  \toprule
  \textbf{Item} & \textbf{Detail} \\
  \midrule
  Process &
    \texttt{python3 apps/fetch-worker/run.py}
    (Docker CMD: \texttt{scripts/fetch\_scheduler\_worker.py}) \\
  Loop &
    \texttt{relocation\_jobs.fetch.scheduler.run\_scheduler\_loop} \\
  Gate &
    \texttt{FETCH\_SCHEDULE\_ENABLED} (code default off) \\
  Interval &
    \texttt{FETCH\_SCHEDULE\_INTERVAL\_HOURS} (default 6, min 0.25) \\
  Countries &
    \texttt{FETCH\_SCHEDULE\_COUNTRIES}, or all
    \texttt{supported\_countries()} \\
  Concurrency &
    \texttt{FETCH\_SCHEDULE\_CONCURRENCY}
    (code default 2, cap 16) \\
  One-shot &
    \texttt{--once} \\
  \bottomrule
\end{tabularx}
\end{center}

Each scheduled pass (\texttt{run\_scheduled\_pass}):

\begin{enumerate}
  \item Listing check --- probe existing open URLs, maybe set
        \texttt{closed\_at}.
  \item Country scrapes --- sequential countries, one at a time.
\end{enumerate}

Skipped if schedule is off, \texttt{httpx} is missing, or another fetch is
running.

Scheduler \texttt{user\_id} is \texttt{resolve\_scheduler\_user\_id()}: first
\texttt{is\_admin = 1} user, else \texttt{PANEL\_ADMIN\_USER} (default
\texttt{admin}). \texttt{fetch\_runs.user\_id} is who started the run, not a
per-user overlay.

\subsection{Panel admin country Fetch}

\begin{itemize}
  \item UI: board Fetch (\texttt{static/js/scrape.js}
        \texttt{startCountryFetch}).
  \item API: \texttt{POST /api/fetch} in \texttt{web/routes/fetch.py}.
  \item Auth: admin. Flag: \texttt{PANEL\_SCRAPE\_ENABLED}.
  \item Body: \texttt{country}, optional \texttt{skip\_filled},
        \texttt{concurrency}, \texttt{ats\_type}. \texttt{country=all} is
        rejected.
\end{itemize}

If scrape is disabled the API returns 503. Header Fetch controls are
admin-only.

\texttt{skip\_filled} is implemented in
\texttt{country\_runner.\_companies\_to\_fetch} but the current board JS does
not send it. The scheduled worker always uses \texttt{skip\_filled=False}.

\subsection{Company-level Fetch jobs}

\begin{itemize}
  \item API: \texttt{POST /api/companies/fetch}.
  \item Auth: login and Full Access (\texttt{plan\_is\_full\_access}).
  \item Flag: \texttt{company\_fetch\_enabled()}.
\end{itemize}

If \texttt{PANEL\_COMPANY\_FETCH\_ENABLED} is set, that value wins; if unset,
it follows \texttt{scrape\_enabled()}. Production panel: scrape off, company
fetch on. On start, \texttt{touch\_company\_fetch\_time} stamps
\texttt{companies.updated}, then \texttt{start\_company\_fetch} runs in a
daemon thread.

\subsection{MCP is not a job fetch}

MCP does not start country or company scrapes. Nearby endpoint
\texttt{POST /api/mcp/positions/<idempotency\_key>/fetch-description} GETs one
job description into \texttt{matching\_jobs.description\_text}. It is gated by
\texttt{company\_fetch\_enabled()} because it does outbound HTTP, but it is
not the fetch-worker pipeline.

\texttt{POST /api/companies} can run ATS detection at insert time. That writes
\texttt{ats\_type} / \texttt{ats\_url}; it does not list jobs.

\subsection{Config flags (names only)}

\begin{center}
\small
\begin{tabularx}{\textwidth}{>{\ttfamily}l X}
  \toprule
  \textrm{\textbf{Variable}} & \textrm{\textbf{Role}} \\
  \midrule
  PANEL\_SCRAPE\_ENABLED &
    Country scrape in the panel. Code default on if unset; production panel 0. \\
  PANEL\_COMPANY\_FETCH\_ENABLED &
    Single-company fetch in the panel. Production panel 1. \\
  FETCH\_SCHEDULE\_ENABLED &
    Worker loop. Code default off; Docker worker 1. \\
  FETCH\_SCHEDULE\_INTERVAL\_HOURS &
    Sleep between passes. Default 6. \\
  FETCH\_SCHEDULE\_CONCURRENCY &
    Companies in flight. Code default 2.
    Dockerfile ENV 4; deploy overrides to 2. \\
  FETCH\_SCHEDULE\_COUNTRIES &
    Optional subset. Empty = all supported countries. \\
  FETCH\_COMPANY\_TIMEOUT\_SECONDS &
    Per-company wait. Default 300. \\
  FETCH\_COUNTRY\_TIMEOUT\_SECONDS &
    Whole country. Default 2700. \\
  PLAYWRIGHT\_BOARD\_TIMEOUT\_SECONDS &
    Playwright listing wall clock. Default 90. \\
  FETCH\_LISTING\_CHECK\_ENABLED &
    Default on. \\
  FETCH\_LISTING\_CHECK\_LIMIT &
    Default 200 open jobs per cycle. \\
  FETCH\_LISTING\_CHECK\_CONCURRENCY &
    Default 2. \\
  FETCH\_LISTING\_CHECK\_MISSES &
    Closed signals before \texttt{closed\_at}. Default 2. \\
  DATABASE\_URL &
    Required Postgres. \\
  SQS\_USER\_OPPORTUNITY\_REFRESH\_QUEUE\_URL
  or ROLE\_PROPAGATOR\_BIN &
    Post-run enqueue; raises if both missing. \\
  \bottomrule
\end{tabularx}
\end{center}

\texttt{supported\_countries()} is \texttt{custom\_countries} labels (seeded
germany / netherlands / uk / portugal, plus panel-added keys such as Ireland,
plus aggregator seeds \texttt{remote-ok}, \texttt{remote-dxb},
\texttt{remote-joblet} created at worker bootstrap). Ireland is not in
\texttt{DEFAULT\_COUNTRY\_LABELS}.

\section{Entry points and process model}

\subsection{Two processes, one catalog}

\begin{lstlisting}
apps/fetch-worker/run.py   scheduled countries + listing check
apps/panel/run.py          country Fetch / company Fetch when flags allow
        |
        v
relocation_jobs/fetch/*    when/what to scrape, run records
relocation_jobs/scrape/*   ATS HTTP / Playwright
relocation_jobs/catalog/repo.py   SQL writes
Postgres                   shared
\end{lstlisting}

Production panel image has no Playwright (\texttt{Dockerfile.ec2}). Worker
image installs Chromium (\texttt{Dockerfile.ec2-worker}). Playwright-only
boards on the slim panel return empty or fail.

\subsection{Country fetch}

\textbf{Worker} (\texttt{run\_country\_fetch\_blocking}): no extra thread. The
scheduler thread calls \texttt{\_country\_fetch\_worker} directly.

\textbf{Panel} (\texttt{start\_country\_fetch}): daemon
\texttt{threading.Thread} so Flask can return \texttt{run\_id}. UI polls
\texttt{GET /api/fetch/status}.

Both eventually:

\begin{lstlisting}
_country_fetch_worker
  asyncio.run(
    async with make_fetch_client(concurrency=workers) as client:
      await asyncio.wait_for(run_country_fetch(...), timeout=country_timeout)
  )
\end{lstlisting}

\texttt{country\_runner.run\_country\_fetch}:

\begin{itemize}
  \item Load stubs: \texttt{list\_country\_company\_stubs}.
  \item Drop \texttt{ats\_type == sourced} (aggregator fan-out employers).
  \item Optional \texttt{ats\_type} filter; optional \texttt{skip\_filled}.
  \item \texttt{asyncio.Semaphore(workers)} + \texttt{asyncio.gather}.
  \item Enrich concurrency = \texttt{max(1, min(4, workers))}.
  \item Shared \texttt{httpx.AsyncClient} (15s timeout,
        \texttt{max\_connections = concurrency + 4}).
\end{itemize}

There is no \texttt{ThreadPoolExecutor} per company. Playwright sync work uses
\texttt{asyncio.to\_thread} and \texttt{\_playwright\_sem = 1}.

\texttt{MAX\_CONCURRENCY} / \texttt{DEFAULT\_CONCURRENCY} are both 16. Panel
country Fetch defaults to 16 unless \texttt{localStorage.panel\_concurrency}
is set. Production worker is 2.

\subsection{Single-company fetch}

\texttt{start\_company\_fetch} $\rightarrow$ daemon thread $\rightarrow$
\texttt{\_company\_fetch\_worker} $\rightarrow$
\texttt{asyncio.run(fetch\_and\_persist\_company(..., review\_mode=True))}.
Uses \texttt{make\_fetch\_client(concurrency=MAX\_CONCURRENCY)} (16). Review
payload is stored on the fetch run for the modal.

\subsection{Playwright (two uses)}

\begin{enumerate}
  \item ATS detection (\texttt{detect\_ats\_via\_playwright}): load careers URL,
        intercept XHR against \texttt{XHR\_ATS\_PATTERNS}, else scan HTML.
  \item Board listing: \texttt{generic} (HTTP HTML then Playwright);
        \texttt{hibob} (XHR intercept then DOM); \texttt{jibe} / \texttt{atlassian}
        (Playwright is primary).
\end{enumerate}

\texttt{fetch\_with\_playwright\_fallback} in \texttt{boards/\_async.py} is
defined but unused. Playwright listing wrap uses
\texttt{ThreadPoolExecutor(max\_workers=1)} so it can time out at
\texttt{PLAYWRIGHT\_BOARD\_TIMEOUT\_SECONDS}; timeout returns an empty list.

\section{Per-company flow}

\texttt{pipeline.fetch\_and\_persist\_company} is the unit of work.

\begin{lstlisting}
get_company(country, name)
start_company_attempt
process_company
  scrape_company_board
    fetch_ats_board
      ensure_company_ats          # detect + persist ats_type/ats_url
      dispatch _BOARD_FETCHERS[ats_type]
    filter_relevant_jobs          # title include/exclude
    filter_jobs_by_expected_locations
    re-attach already-known URLs from raw
    merge_matching_jobs
    enrich_jobs (only_missing=True)
    mark fetch_ok
    sync_company_board_to_catalog
finish_company_attempt
\end{lstlisting}

\subsection{First run vs later runs}

\texttt{ensure\_company\_ats}:

\begin{enumerate}
  \item \texttt{apply\_known\_ats\_override} --- \texttt{KNOWN\_ATS} /
        \texttt{FORCE\_KNOWN\_ATS} and URL detectors.
  \item If \texttt{ats\_type} is \texttt{generic}: stop. Generic is sticky.
  \item If any other non-empty type: use cached \texttt{ats\_url}.
  \item If empty: \texttt{KNOWN\_ATS} else \texttt{detect\_ats\_for\_company}
        (static HTML, else Playwright).
  \item Persist via \texttt{sync\_board} \textbf{before} the board API call.
\end{enumerate}

\texttt{add\_company} already detects ATS, so a newly added company often has
\texttt{ats\_type} set before the first job fetch.

\subsection{Aggregator companies}

If \texttt{ats\_type} is \texttt{remoteok} / \texttt{remotedxb} /
\texttt{joblet}: fetch the feed, group by employer,
\texttt{sync\_aggregator\_employer\_jobs} upserts \texttt{sourced} employer
companies with additive job inserts. The aggregator company's own
\texttt{matching\_jobs} is cleared. Country fetch skips \texttt{sourced}
stubs.

\section{ATS types}

Display/API list: \texttt{ATS\_TYPE\_CHOICES} in
\texttt{core/ats\_constants.py}. Runtime dispatch:
\texttt{\_BOARD\_FETCHERS} in \texttt{scrape/board.py}. \texttt{generic} is
not in \texttt{ATS\_TYPE\_CHOICES}. Unknown ids raise
\texttt{UnsupportedAtsTypeError}.

README says ``25 ATS types''. Code has \textbf{30} named types plus generic:

ashby, atlassian, applytojob, bamboohr, bol, deel, epam, greenhouse,
greenhouse\_eu, hirehive, jibe, job\_shop, joblet, join, lever, lever\_eu,
movingimage, personio, project\_a, recruitee, remotedxb, remoteok, rss,
smartrecruiters, teamtailor, workable, workday, hibob, pinpointhq,
successfactors.

Playwright-primary: atlassian, jibe. HiBob intercepts \texttt{/api/job-ad},
then DOM fallback. Greenhouse uses
\texttt{boards-api.greenhouse.io/v1/boards/\{slug\}/jobs}. Workday POSTs CXS
\texttt{/jobs}.

Detection order: optional user hint, \texttt{KNOWN\_ATS}, static URL/HTML,
Playwright XHR, else generic.

\section{Relevance filtering}

Applied after the raw board list, before merge
(\texttt{scrape/company.\_filter\_board\_listings}):

\begin{enumerate}
  \item \texttt{filter\_relevant\_jobs(raw, True)} $\rightarrow$
        \texttt{is\_relevant(title)}.
  \item \texttt{filter\_jobs\_by\_expected\_locations} --- office-tag gate.
\end{enumerate}

\texttt{relevant\_only=True} is hard-coded. Keywords live in
\texttt{core/ats\_constants.py}; logic in \texttt{scrape/relevance.py}.

Include (title must contain at least one): backend, software engineer,
platform engineer, golang / go engineer, java / javascript / typescript,
kotlin, python engineer, fullstack, product engineer, senior engineer, and
the rest of \texttt{INCLUDE\_KEYWORDS}.

Exclude plus extra rules: CTO, staff unless senior/staff, cloud engineer or
AI platform without backend/software. If the title has
engineer/developer/programmer, marketing and hr are not applied as exclude
keywords. Jobs without title or URL are dropped.

Single-company fetch records why titles were dropped
(\texttt{scrape/review.explain\_title\_filter}) into
\texttt{fetch\_runs.review\_jobs\_json}.

Location gate: listings with no usable location text are kept; remote-only
are excluded; explicit other country/city excluded. Already-catalogued URLs
are re-attached from the raw list even if they fail title/location this run,
so they stay open and get location fields updated.

On panel read, known-wrong location can still be routed to Not for me
(\texttt{panel/flatten\_jobs.py}). That is overlay/read behavior, not a fetch
write.

\section{Data model}

Schema is created at \texttt{init\_db()} plus
\texttt{catalog/schema.init\_catalog\_schema()} and
\texttt{core/migrations.py}. There is no separate Alembic tree.

\subsection{Catalog (shared --- fetch writes here)}

\textbf{\texttt{country\_meta}}: \texttt{country} PK, \texttt{source},
\texttt{fetched}, \texttt{updated}, \texttt{jobs\_fetched}, \texttt{total},
\texttt{last\_fetch\_new\_jobs}.

\textbf{\texttt{companies}}: unique \texttt{(country, name)}; \texttt{city},
\texttt{size}, \texttt{careers\_url}, \texttt{ats\_type}, \texttt{ats\_url}
(ATS cache), \texttt{fetch\_problem} / dates, \texttt{fetch\_ok} / dates,
\texttt{sources\_json}, \texttt{cities\_json}, \texttt{locations\_json},
\texttt{catalog\_kind} (\texttt{relocation} or \texttt{remote}).

\textbf{\texttt{matching\_jobs}}: FK \texttt{company\_id} ON DELETE CASCADE;
unique \texttt{(company\_id, idempotency\_key)}; \texttt{title}, \texttt{url},
\texttt{fetched}, \texttt{last\_seen}, \texttt{visa\_sponsorship},
\texttt{location}, \texttt{locations\_json}, \texttt{description\_text},
\texttt{public\_slug}, \texttt{closed\_at}, \texttt{listing\_misses}.

There is no orphans table. Orphan is a read-time concept: a tracking row whose
job is closed or missing from the live board.

\subsection{Operational (fetch writes here)}

\textbf{\texttt{fetch\_runs}}: \texttt{user\_id}, \texttt{country},
\texttt{company\_name}, \texttt{scope} (country or company), \texttt{status}
(running / done / failed / cancelled), timestamps, \texttt{exit\_code},
cancel flags, \texttt{new\_jobs}, \texttt{concurrency}, progress/activity/log
JSON (UI log capped at 200 lines), \texttt{review\_jobs\_json}.
\texttt{file\_name} (e.g.\ \texttt{uk\_companies.json}) is a legacy archive
label; catalog persistence is Postgres.

\textbf{\texttt{company\_fetch\_attempts}}: \texttt{fetch\_run\_id},
\texttt{country}, \texttt{company\_name}, \texttt{careers\_url},
\texttt{ats\_type}, times, \texttt{status}, \texttt{error\_message},
\texttt{jobs\_total}, \texttt{jobs\_new}, \texttt{message},
\texttt{duration\_seconds}.

\subsection{Per-user overlay (fetch does not write)}

\texttt{job\_tracking}, \texttt{company\_tracking}, \texttt{job\_status\_events}.
\texttt{user\_opportunities} / \texttt{position\_broadcast\_assignments} are
written by the role-propagator after enqueue. \texttt{positions/} is the user
mutation service; fetch never calls it.

Board \texttt{GET /api/board} merges catalog with tracking in
\texttt{panel/flatten\_jobs.py}. Closed unengaged roles are hidden
(\texttt{skip\_closed\_unengaged}). If the user still has tracking and the role
is gone, \texttt{flatten\_orphans.append\_tracked\_orphans} can reinject it.

\section{Merge semantics}

In-memory merge: \texttt{scrape/merge.merge\_matching\_jobs}. Identity:
\texttt{job\_idempotency\_key(url)} --- SHA-256 of \texttt{job:v1:} plus
normalized URL (lowercase host, strip www, keep only job-id query params,
drop fragment, trim slash). README ``merge by URL'' is this key.

\begin{itemize}
  \item Existing key on the board: keep original \texttt{fetched}, update
        title/url/description/location, copy visa from old if set, clear
        \texttt{closed\_at} and \texttt{listing\_misses}, set
        \texttt{last\_seen} now (preserved).
  \item New key: insert with \texttt{fetched}/\texttt{last\_seen} now.
  \item Existing key not in this scrape: keep the old dict, set
        \texttt{closed\_at} if empty (stale kept).
\end{itemize}

Scrape does not delete roles in the merge step. Persistence upserts the full
merged list and \texttt{DELETE}s rows whose keys are not in that list. Because
stale jobs remain in the merged list, they stay in Postgres with
\texttt{closed\_at}. If the same job returns later, merge clears
\texttt{closed\_at}. \texttt{public\_slug} is preserved on update.

Visa text is filled later by \texttt{enrich\_jobs}
(\texttt{detect\_visa\_relocation} on JD text), \texttt{only\_missing=True}.

\section{Write path}

SQL is confined to \texttt{*/repo.py} and schema/migrations.

\texttt{sync\_company\_board\_to\_catalog} (one \texttt{db\_transaction}):

\begin{enumerate}
  \item Upsert company on \texttt{(country, name)} (does not overwrite
        \texttt{added}).
  \item Per-job insert on conflict \texttt{(company\_id, idempotency\_key)},
        then delete keys not in the merged list.
  \item Patch \texttt{country\_meta.updated}.
\end{enumerate}

Attempts and live run progress are separate short transactions. Company A can
commit while B is still scraping.

Post-run enqueue is after \texttt{persist\_fetch\_run}. If SQS and
\texttt{ROLE\_PROPAGATOR\_BIN} are both unset, it raises; catalog is already
committed.

\section{Errors, retries, logging}

Per company, non-cancel exceptions: infra messages (can't start new thread,
cannot allocate memory, too many open files) do not set
\texttt{fetch\_problem}. Other errors mark \texttt{fetch\_problem}, still
\texttt{sync\_board}, attempt status error, \texttt{jobs\_new=0}.

Timeouts: company 300s --- log and continue the country; country 2700s --- run
failed; Playwright board timeout --- empty list. No automatic retry of a
failed company inside the cycle.

Listing check: two CLOSED probes set \texttt{closed\_at}; 429/403/5xx/timeout
are UNKNOWN and do not increment misses.

Cancel: \texttt{POST /api/fetch/cancel} sets \texttt{cancel\_requested}.
Country runner polls at most once per second. Exit code 130.

Orphan reap marks all \texttt{status=running} rows failed. Called from worker
bootstrap and from \texttt{reap\_zombie\_fetch} when this process has no live
local thread and \texttt{fetch\_process\_may\_reap\_orphans()} is true (scrape
on or company-fetch on). Production panel has company fetch on, so status
polls can reap a worker's running row. Ops doc describes a narrower intended
rule; the code is broader.

Logs: structlog logger \texttt{relocation\_jobs.fetch}; JSON if stderr is not
a TTY; UI buffer last 200 lines; Grafana/Loki on EC2. Docker logs wipe on
deploy. Durable history is Postgres attempts/runs.

\section{Production vs local}

\begin{center}
\small
\begin{tabularx}{\textwidth}{l X X X}
  \toprule
  & \textbf{Local panel} & \textbf{EC2 panel} & \textbf{EC2 worker} \\
  \midrule
  Image &
    \texttt{apps/panel/run.py} &
    \texttt{Dockerfile.ec2} &
    \texttt{Dockerfile.ec2-worker} \\
  Playwright &
    expected &
    not installed &
    Chromium headless shell \\
  PANEL\_SCRAPE\_ENABLED &
    1 &
    0 &
    1 \\
  Company fetch &
    yes &
    yes (httpx ATS) &
    not via HTTP \\
  Schedule &
    off unless worker &
    unused &
    every 6h, concurrency 2 \\
  Country Fetch &
    admin, in-process &
    disabled &
    scheduler \\
  Listing check &
    only if worker &
    no &
    yes, before countries \\
  \bottomrule
\end{tabularx}
\end{center}

Deploy worker env includes schedule 6h, concurrency 2, timeouts,
\texttt{DATABASE\_URL}, SQS. It does not set
\texttt{FETCH\_SCHEDULE\_COUNTRIES} (all supported, including remote
aggregators after seed). Slim panel company fetch cannot run Playwright
detection/boards; Greenhouse/Lever/Ashby-style APIs work in-process.

\section{End-to-end happy path}

Assume the worker is up, catalog has UK companies with cached Greenhouse
URLs, and no other fetch is running.

\begin{enumerate}
  \item \texttt{run\_scheduler\_loop} wakes (or \texttt{--once}).
  \item Listing check: up to 200 open jobs, maybe
        \texttt{UPDATE matching\_jobs.listing\_misses/closed\_at}.
  \item Sequential countries. For \texttt{uk}: insert
        \texttt{fetch\_runs} (\texttt{status=running}, \texttt{scope=country}).
  \item One asyncio loop, one httpx client, semaphore 2.
  \item Stubs: names + \texttt{ats\_type} + \texttt{has\_jobs}; drop sourced.
  \item Company Acme: load jobs, insert attempt, cached greenhouse, GET
        boards-api jobs.
  \item Title + location filter. Merge: new URLs added, old URLs keep
        \texttt{fetched}, missing URLs get \texttt{closed\_at}.
  \item Enrich JDs only where visa/description/location incomplete.
  \item Upsert company (\texttt{fetch\_ok=1}) and jobs in one transaction.
  \item Attempt status ok. After all UK companies: patch
        \texttt{country\_meta}. Finalize run \texttt{status=done}.
  \item If success and companies done $>$ 0: SQS
        \texttt{\{type: country, country: uk\}}.
  \item Sleep 6 hours between full passes (not between countries).
\end{enumerate}

\begin{lstlisting}
scheduler
  -> listing_check
  -> INSERT fetch_runs running
  -> asyncio country loop (semaphore)
       -> ATS API / Playwright
       -> filter / merge / enrich
       -> upsert companies + matching_jobs
  -> fetch_runs status=done
  -> SQS country refresh
\end{lstlisting}

\section{Key file map}

\begin{center}
\small
\begin{tabularx}{\textwidth}{>{\ttfamily}l X}
  \toprule
  \textrm{\textbf{Path}} & \textrm{\textbf{Role}} \\
  \midrule
  apps/fetch-worker/run.py &
    Worker entry \\
  scripts/fetch\_scheduler\_worker.py &
    Docker/ops shim \\
  Dockerfile.ec2-worker &
    Playwright worker image \\
  Dockerfile.ec2 &
    Slim panel; scrape off \\
  scripts/build\_companies.py &
    Discovery only --- not job fetch \\
  core/panel\_flags.py &
    Scrape / company-fetch / orphan-reap gates \\
  core/ats\_constants.py &
    Caps, keywords, ATS choices, KNOWN\_ATS \\
  core/ats\_detection.py &
    Static + Playwright detection \\
  core/job\_identity.py &
    Normalize URL $\rightarrow$ SHA-256 key \\
  fetch/scheduler.py &
    Interval loop, listing check then fetch \\
  fetch/runner.py &
    Threads vs blocking; enqueue refresh \\
  fetch/country\_runner.py &
    Semaphore fan-out \\
  fetch/pipeline.py &
    Load $\rightarrow$ scrape $\rightarrow$ persist \\
  fetch/service.py &
    Attempt lifecycle; fetch\_problem vs infra \\
  fetch/repo.py &
    SQL for runs and attempts \\
  fetch/listing\_check.py &
    Pre-scrape URL closer \\
  scrape/company.py &
    Filter, merge, enrich \\
  scrape/board.py &
    ATS dispatcher \\
  scrape/ats\_resolve.py &
    First-run detect + cache \\
  scrape/merge.py &
    Merge by idempotency key \\
  scrape/relevance.py &
    Title include/exclude \\
  scrape/enrich.py &
    JD + visa flag \\
  catalog/repo.py &
    \texttt{sync\_company\_board\_to\_catalog} \\
  web/routes/fetch.py &
    Country Fetch / status / cancel \\
  web/routes/companies.py &
    Company Fetch (Full Access) \\
  panel/flatten\_jobs.py &
    Read merge; hide closed unengaged \\
  panel/flatten\_orphans.py &
    Reinject tracked missing jobs \\
  positions/service.py &
    User tracking writes (not used by fetch) \\
  async\_jobs/enqueue.py &
    Post-fetch SQS / bin enqueue \\
  \bottomrule
\end{tabularx}
\end{center}

\section{Doc vs code notes}

\begin{itemize}
  \item ATS count: README 25; \texttt{ATS\_TYPE\_CHOICES} has 30 plus generic.
  \item Concurrency: README ``up to 16 workers'' is the cap / local panel
        default. Production worker is 2 via deploy, even though
        \texttt{Dockerfile.ec2-worker} ENV is 4.
  \item Ireland: marketed as a country; not in seed defaults; custom
        \texttt{custom\_countries} row.
  \item Orphan reap: ops doc says the panel should not reap worker runs unless
        scrape is on; code also treats
        \texttt{PANEL\_COMPANY\_FETCH\_ENABLED=1} as permission to reap all
        running rows.
  \item \texttt{skip\_filled}: API supports it; current Fetch UI does not send
        it.
  \item MCP does not trigger fetch-worker jobs; only optional JD text fetch.
\end{itemize}

\section{Invariants}

\begin{enumerate}
  \item Companies first, jobs second. Empty country catalog cannot fetch.
  \item Postgres is the catalog. No JSON job pipeline.
  \item Scrape does not delete roles. Missing from ATS sets
        \texttt{closed\_at}.
  \item One fetch at a time. Countries sequential; companies concurrent.
  \item SQL stays in repos.
  \item Successful \textbf{run} enqueues country refresh; scrape code does not.
  \item Infra thread/memory/file errors do not stamp
        \texttt{fetch\_problem}.
  \item Fetch does not write per-user tracking.
\end{enumerate}

\end{document}
